% Generated draft: requires mathematical and source review.
\documentclass[11pt]{article}
\usepackage[margin=25mm]{geometry}
\usepackage{fontspec}
\setmainfont{Latin Modern Roman}
\usepackage{amsmath,amssymb,mathtools,unicode-math}
\setmathfont{Latin Modern Math}
\usepackage{enumitem,needspace,xurl,hyperref}
\setlength{\parindent}{0pt}
\setlength{\parskip}{5pt}
\setlength{\emergencystretch}{3em}
\newcommand{\N}{\mathbb{N}}
\newcommand{\Z}{\mathbb{Z}}
\newcommand{\Q}{\mathbb{Q}}
\newcommand{\R}{\mathbb{R}}
\newcommand{\C}{\mathbb{C}}
\newcommand{\F}{\mathbb{F}}
\newcommand{\A}{\mathbb{A}}
\newcommand{\PP}{\mathbb{P}}
\newcommand{\E}{\mathbb{E}}
\newcommand{\eps}{\varepsilon}
\newcommand{\dd}{\,\mathrm{d}}
\begin{document}
\section*{UnsolvedMath ID: 800422}
\textbf{Source problem code:} catalogue entry
\subsection{Definitions and mathematical statement}
Let \(K\) be a number field, and let \(f\in K(x)\) have degree \(d\geq 2\). Regard \(f\) as a morphism \(\PP^1\to\PP^1\), so that iteration is defined even when an orbit passes through infinity. Write \(f^{\circ n}\) for its \(n\)-fold iterate. A point \(\alpha\in K\subset\PP^1(K)\) is periodic if \(f^{\circ m}(\alpha)=\alpha\) for some integer \(m\geq 1\).

For each integer \(n\geq 1\), choose relatively prime polynomials \(P_n,Q_n\in K[x]\) such that
\[
f^{\circ n}(x)=\frac{P_n(x)}{Q_n(x)}.
\]
Factor the nonzero polynomial \(P_n-\alpha Q_n\) over \(K\) as
\[
P_n(x)-\alpha Q_n(x)=c_n\prod_{j=1}^{r_n}h_{n,j}(x)^{e_{n,j}},
\]
where \(c_n\in K^\times\), the \(h_{n,j}\) are distinct monic irreducible polynomials of positive degree, and \(e_{n,j}\geq 1\). For a nonzero constant polynomial the product is empty. Use the following interpretation of eventual stability over \(K\):
\[
(f,\alpha)\text{ is eventually stable over }K
\quad\Longleftrightarrow\quad
\exists B\in\Z_{\geq 0}\ \forall n\geq 1,
\qquad \sum_{j=1}^{r_n}e_{n,j}\leq B.
\]
Thus irreducible factors are counted with multiplicity. The supplied record does not define eventual stability; this explicit conventional interpretation, including the multiplicity convention, requires confirmation against the cited source.

\textbf{Conjecture (Jones--Levy, 2017; supplied reference MR3704363, Conjecture 1.2).} For every number field \(K\), every rational function \(f\in K(x)\) of degree at least \(2\), and every \(\alpha\in K\) satisfying
\[
\forall m\geq 1,\qquad f^{\circ m}(\alpha)\neq\alpha,
\]
the pair \((f,\alpha)\) is eventually stable over \(K\). The bound \(B\) may depend on \(K\), \(f\), and \(\alpha\). The hypothesis excludes periodic points but does not exclude nonperiodic preperiodic points.
\subsection{Short English statement}
For a rational map of degree at least two over a number field and a nonperiodic \(K\)-rational point \(\alpha\), the conjecture predicts a uniform bound on the number of irreducible factors over \(K\) of the numerators of \(f^{\circ n}(x)-\alpha\), counted with multiplicity, as \(n\) varies. The supplied registry labels this conjecture open; that status is not independently verified.
\subsection{Sources}
\begin{itemize}
\item Supplied local UnsolvedMath registry record, permanent ID \(800422\), exact source problem code \texttt{catalogue entry}; upstream revision \texttt{d371ee917ac9647fbb06cd6fc529deec6f3dda82}. The record attributes the source list to Benedetto, DeMarco, Ingram, Jones et al., with proposed year 2018, and supplies the status \texttt{open}.
\item Benedetto, DeMarco, Ingram, Jones et al., \emph{Current Trends and Open Problems in Arithmetic Dynamics} (2018), Conjecture 45. Supplied source URL: \url{https://arxiv.org/abs/1806.04980}. The record reports access on 2026-10-16 and extraction by \texttt{arxiv\_latex}; these are supplied provenance details, not independent source verification.
\item The supplied conjecture text attributes the assertion to Jones--Levy (2017), reference MR3704363, Conjecture 1.2. Further bibliographic details for that reference are not supplied.
\end{itemize}
\end{document}
